% !PRINTED 609
\weberpage{609}
In the same case are found the cyclic and in general all irreducible
\gesperrt{Abel}ian equations of odd degree, which therefore can never be solved by real
radicals.

Accordingly we can, for example, distinguish the cubic equations in the field of the
rational numbers into three or four kinds, which all occur in certain geometric
problems famous from antiquity. We here leave the reducible equations out of account.
The degree of the \gesperrt{Galois} group must then always be divisible by 3
(\S.\,154, 7.) and is therefore either equal to 3 or equal to 6. By \S.\,157 the group
can therefore never be reduced by adjunction of mere square roots to a lower than
the third degree; hence the equation also cannot be solved by
square roots. The corresponding geometric problems cannot be solved with
compass and ruler.

If the group is of degree 3, we have a cyclic equation. Here belong the
equations arising from the division of the circle, e.g. the one on which the construction
of the regular heptagon depends. The three roots of such an equation are real
and cannot be expressed by real radicals. Among the cubic
equations with a group of degree $6$ we must distinguish those with
positive and those with negative discriminant.

The first belong to the casus irreducibilis and can be reduced to the
equation on which the trisection of an arbitrary angle depends. To the
equations with negative discriminant belong also, as special cases, the pure
cubic equations $x^3 = a$, when $a$ is not a cube number. For $a = 2$ there results
the equation on which the \gesperrt{Delian problem} of the duplication of the cube depends.

{\centering\S.\,180.\par}
{\centering \gesperrt{Metacyclic equations of prime degree.}\par}

The general conditions which we found in \S.\,176 are not yet
simple enough to permit an immediate application to the finding of metacyclic
equations or to the decision on the solvability of a proposed
%CONT>>
